\documentclass[a4paper]{article}

% if you need to pass options to natbib, use, e.g.:
%     \PassOptionsToPackage{numbers, compress}{natbib}
% before loading neurips_2022


% ready for submission
%\usepackage{ofu_xai_2022}

\input{math_commands.tex}

% to compile a preprint version, e.g., for submission to arXiv, add add the
% [preprint] option:
%     \usepackage[preprint]{ofu_xai_2022}


% to compile a camera-ready version, add the [final] option, e.g.:
\usepackage[final, nonatbib]{ofu_xai_2022}

% to avoid loading the natbib package, add option nonatbib:
%    \usepackage[nonatbib]{ofu_xai_2022}


\usepackage[utf8]{inputenc} % allow utf-8 input
\usepackage[T1]{fontenc}    % use 8-bit T1 fonts
\usepackage{hyperref}       % hyperlinks
\usepackage{url}            % simple URL typesetting
\usepackage{booktabs}       % professional-quality tables
\usepackage{amsfonts}       % blackboard math symbols
\usepackage{nicefrac}       % compact symbols for 1/2, etc.
\usepackage{microtype}      % microtypography
\usepackage{xcolor}         % colors
\usepackage{graphicx}


\usepackage[round]{natbib}
%%%%

\title{Evaluating CNN Architectures and Data Augmentation for Object Recognition}


% The \author macro works with any number of authors. There are two commands
% used to separate the names and addresses of multiple authors: \And and \AND.
%
% Using \And between authors leaves it to LaTeX to determine where to break the
% lines. Using \AND forces a line break at that point. So, if LaTeX puts 3 of 4
% authors names on the first line, and the last on the second line, try using
% \AND instead of \And before the third author name.


\author{%
  Elia \\
  Otto-Friedrich University of Bamberg\\
  96049 Bamberg, Germany\\
  \texttt{max.mustermann@stud.uni-bamberg.de} \\[0.5cm]
  Seminar: xAI-Sem-?? \\
  Degree: M.Sc. (CitH/AI/\dots) \\
  Matriculation \#: 12345678
  % examples of more authors  
  \And
  Dominic \\
  Otto-Friedrich University of Bamberg\\
  96049 Bamberg, Germany\\
  \texttt{max.mustermann@stud.uni-bamberg.de} \\[0.5cm]
  Seminar: xAI-Sem-?? \\
  Degree: M.Sc. (CitH/AI/\dots) \\
  Matriculation \#: 12345678
  \And
  Magda \\
  Otto-Friedrich University of Bamberg\\
  96049 Bamberg, Germany\\
  \texttt{magdalena\_klug@stud.uni-bamberg.de} \\[0.5cm]
  Proejct: xAI-Proj-B \\
  Degree: B.Sc. AI \\
  Matriculation \#: 12345678
  % Coauthor \\
  % Affiliation \\
  % Address \\
  % \texttt{email} \\
  % \AND
  % Coauthor \\
  % Affiliation \\
  % Address \\
  % \texttt{email} \\
  % \And
  % Coauthor \\
  % Affiliation \\
  % Address \\
  % \texttt{email} \\
  % \And
  % Coauthor \\
  % Affiliation \\
  % Address \\
  % \texttt{email} \\
}


\begin{document}


\maketitle
\def\va{{\bm{a}}}

\begin{abstract}
  The abstract paragraph should be indented \nicefrac{1}{2}~inch on
  both the left- and right-hand margins. Use 10~point type, with a vertical
  spacing (leading) of 11~points.  The word \textbf{Abstract} must be centered,
  bold, and in point size 12. Two line spaces precede the abstract. The abstract
  must be limited to one paragraph.
\end{abstract}


\section{Introduction (E)}


Please read the instructions below carefully and follow them faithfully. Those instructions and style file is based on \verb+neurips_2022.tex+ (LPPL v1.3c) and was adapted for the purposes at the Chair of Explainable Machine Learning at the University of Bamberg.

\section{Methodology}

\subsection{Data Augmentation (E)}
\subsection{Architectures (E)}
\begin{itemize}
\item SimpleNet
\item LargerNet
\item ResNet
\end{itemize}
\subsection{Training (M)}
Model training was conducted using the Adam optimizer \citep{kingma2017} with default settings (initial learning rate of 0.0001). To ensure convergence and prevent overfitting, we employed a learning rate scheduler and an early stopping mechanism. 

The learning rate scheduler implemented a ReduceLROnPlateau strategy, which reduces the learning rate by a factor of 0.5 when the validation loss fails to decrease over a patience period of 2 epochs. This adaptive learning rate adjustment enables finer gradient updates as the model approaches convergence.

Early stopping was applied to halt training when model performance on the validation set ceased to improve. Specifically, training was terminated if the validation loss did not decrease by a minimum delta of 0.1 over a patience period of 5 epochs. This approach prevents overfitting while minimizing computational overhead \citep{prechelt2002}.

The training procedure followed an iterative process over multiple epochs: each epoch consisted of a training phase followed by validation. Subsequently, performance metrics were computed and logged, the learning rate scheduler was invoked to adjust the learning rate if necessary, model checkpoints were saved if the current validation metrics have been better than previous epochs and the early stopping criterion was evaluated to determine training continuation. If early stopping hasn't been triggered, the training stops after a total of 300 epochs.
\subsection{Evaluation Metrics (M)}
Model performance was assessed using multiple evaluation metrics to provide a comprehensive analysis of classification quality.
\subsubsection{Accuracy}
Accuracy was calculated as the percentage of correct predictions relative to the total number of predictions. For binary or multi-class classification, accuracy is defined as:

\begin{equation}
\text{Accuracy} = \frac{TP + TN}{TP + TN + FP + FN}
\end{equation}

where $TP$, $TN$, $FP$, and $FN$ represent true positives, true negatives, false positives, and false negatives, respectively. In the context of this study, the confusion matrix extends to $10 \times 10$ dimensions, accommodating all class predictions.

\subsubsection{ROC and AUC}
Receiver Operating Characteristic (ROC) curve analysis and the corresponding Area Under the Curve (AUC) were employed to evaluate the model's ability to distinguish between classes \citep{hand2001}. For multi-class classification, we adopted a One-vs-Rest (OvR) approach following \cite{fawcett2006}, calculating individual ROC curves for each class against all other classes. The True Positive Rate (TPR) and False Positive Rate (FPR) for each class are computed as follows:

\begin{equation}
\text{TPR} = \frac{TP}{TP + FN}
\end{equation}

\begin{equation}
\text{FPR} = \frac{FP}{FP + TN}
\end{equation}

The overall model performance was quantified using the weighted-averaged AUC, which computes the weighted mean of individual class AUC values, thereby treating all classes regarding their sample size. This approach provides a balanced assessment of model performance across all classes.

\section{Experiments (M)}
\subsection{Training Setup}
%TODO add hard facts: kaggle, gpu/cpu, batch size, framework (pytorch mit version), python version, trainingszeit, random seeds, gewichtsinitalisierung

\subsection{Architecture Comparison}
Table~\ref{tab:training_results} summarizes the training results for each of the implemented neural network architectures.

SimpleCNN demonstrated the lowest performance across all metrics. The non-augmented variant stopped training early after only 8 epochs and showed a larger validation-training loss difference compared to the augmented version.

LargerNet achieved intermediate performance levels between SimpleCNN and ResNet. The augmented variant completed the full training cycle of 300 epochs with a smaller loss difference, while the non-augmented version stopped after 91 epochs. Both variants achieved similar accuracy and AUC/ROC scores.

ResNet exhibited the highest performance among all tested architectures, with both variants completing 300 epochs of training. The augmented and non-augmented models achieved comparable accuracy and AUC/ROC values, with the augmented variant showing a smaller validation-training loss difference.

Across all architectures, data augmentation consistently reduced the gap between validation and training loss, with the most substantial effect observed in SimpleCNN.


\begin{table}[h]
\centering
\caption{Training results on validation data for different neural network architectures with and without data augmentation.}
\label{tab:training_results}
\begin{tabular}{lcccccc}
\hline
\textbf{Model} & \multicolumn{2}{c}{\textbf{SimpleCNN}} & \multicolumn{2}{c}{\textbf{LargerNet}} & \multicolumn{2}{c}{\textbf{ResNet}} \\
\cline{2-7}
\textbf{Augmentation} & Yes & No & Yes & No & Yes & No \\
\hline
Epochs trained & 30 & 8 & 300 & 91 & 300 & 300 \\
Learning Rate & 0.0001 & 0.00005 & $1.5 \times 10^{-8}$ & $1.5 \times 10^{-8}$ & $1.5 \times 10^{-8}$ & $1.5 \times 10^{-8}$ \\
Loss Difference & 1.154 & 6.794 & 0.102 & 0.261 & 0.069 & 0.203 \\
(Val. - Train.) & & & & & & \\
Accuracy & 0.367 & 0.279 & 0.447 & 0.476 & 0.893 & 0.879 \\
AUC / ROC & 0.785 & 0.683 & 0.838 & 0.848 & 0.994 & 0.992 \\
\hline
\end{tabular}
\end{table}


\section{Results}

\subsection{Results on Testset (D)}
All results in this section are based solely on the self-collected outdoor test dataset. This dataset was not used during training or validation and differs from the ImageNet-based training data in terms of environment and image conditions.

The self-collected test dataset consists of 949 outdoor images covering the same ten object classes as the training data. All images were captured by four independent student groups in real-world outdoor environments. The dataset includes variations in background, lighting, and viewpoints and was used exclusively to evaluate model performance under domain shift.

The evaluation aims to analyze how well the trained models generalize to unseen outdoor data. Performance is reported using overall accuracy, confusion matrices, per-class accuracy, and ROC-AUC scores.

\subsection{Results on Testset (D)}
Performance on the self-collected outdoor test dataset is compared across all evaluated models. 

Table~\ref{tab:test_results} reports the overall classification accuracy and ROC-AUC scores for SimpleNet, LargerNet and ResNet18. Each model is evaluated with and without data augmentation.

Across all evaluation settings, ResNet18 achieves the highest performance on the outdoor test set in terms of both accuracy and ROC-AUC. The two custom architectures, SimpleNet and LargerNet show lower performance on the same data. LargerNet consistently outperforms SimpleNet, while both remain clearly below the performance level of ResNet18. 

The relative performance ordering of the evaluated architectures remains consistent across both evaluation metrics.

\begin{table}[h]
\centering
\caption{Test results on the self-collected outdoor dataset for different neural network architectures with and without data augmentation.}
\label{tab:test_results}

\setlength{\tabcolsep}{8pt}           % horizontal padding 
\renewcommand{\arraystretch}{1.15}    % vertical padding

\begin{tabular}{lcccccc}
\hline
\textbf{Metric} & \multicolumn{2}{c}{\textbf{SimpleNet}} & \multicolumn{2}{c}{\textbf{LargerNet}} & \multicolumn{2}{c}{\textbf{ResNet18}} \\
\cline{2-7}
\textbf{Augmentation} & Yes & No & Yes & No & Yes & No \\
\hline
Accuracy      & 0.125 & 0.154 & 0.162 & 0.203 & 0.466 & 0.685 \\[2pt]
ROC-AUC (OvR) & 0.566 & 0.608 & 0.600 & 0.629 & 0.886 & 0.955 \\
\hline
\end{tabular}
\end{table}





\subsection{Effects of Data Augmentation (D)}
The effect of data augmentation on test performance is evaluated by comparing augmented and non-augmented variants of all models. The corresponding accuracy and ROC-AUC values are reported in Table~\ref{tab:test_results}.

For all three architectures, data augmentation results in lower classification accuracy compared to the non-augmented variants. A similar decrease can be observed for the ROC-AUC scores. The largest absolute performance drop occurs for ResNet18, while SimpleNet and LargerNet show smaller differences between the two configurations.

\subsection{Class Separability (D)}

Class separability is analyzed using the best-performing model, ResNet18 without data augmentation.

Figure~\ref{fig:per_class_accuracy_resnet18} shows the per-class classification accuracy on the outdoor test dataset.
The figure indicates noticeable differences in classification performance across object categories.
While classes such as computer-keyboard and mouse achieve comparatively high accuracy, other classes such as binder and notebook show lower accuracy values, suggesting that these classes are more difficult to distinguish in outdoor settings.

\begin{figure}[h]
\centering
\includegraphics[width=0.9\linewidth]{figures/resnet18_per_class_accuracy.png}
\caption{Per-class classification accuracy for ResNet18 without data augmentation on the outdoor test dataset.}
\label{fig:per_class_accuracy_resnet18}
\end{figure}


To analyze the underlying error patterns in more detail, Figure~\ref{fig:confusion_matrix_resnet18} presents the corresponding confusion matrix.
The matrix shows that most predictions are located on the diagonal, indicating correct classification for a large portion of the test samples.
However, several off-diagonal entries reveal recurring confusions between specific class pairs.

\begin{figure}[h]
\centering
\includegraphics[width=0.9\linewidth]{figures/resnet18_confusion_matrix.png}
\caption{Confusion matrix for ResNet18 without data augmentation evaluated on the test dataset.}
\label{fig:confusion_matrix_resnet18}
\end{figure}


In particular, images of teapots are frequently misclassified as coffee mugs, with 22 such cases visible in the confusion matrix.
Similarly, notebooks are misclassified as binders in 12 instances, while binders are predicted as notebooks in 13 instances.
Another prominent confusion occurs between computer keyboards and remote controls, where 17 keyboard images are classified as remote controls.
These confusion patterns correspond to the lower per-class accuracy values observed for the affected classes in Figure~\ref{fig:per_class_accuracy_resnet18}.

\section{Discussion}
\subsection{Discoveries (M)}
\subsection{Limitations (M)}

\section{Conclusion (E)}


\newpage

\bibliography{bibliography}
\bibliographystyle{abbrvnat}


%%%%%%%%%%%%%%%%%%%%%%%%%%%%%%%%%%%%%%%%%%%%%%%%%%%%%%%%%%%%%%%%%%%%%%%%%%%%%%%%%%%%%%%%%%%%%%%%%%%%
%% Declaration of Authorship
%%%%%%%%%%%%%%%%%%%%%%%%%%%%%%%%%%%%%%%%%%%%%%%%%%%%%%%%%%%%%%%%%%%%%%%%%%%%%%%%%%%%%%%%%%%%%%%%%%%%


{\parindent 0cm
%%%%%%%%%%%%%%%%%%%%%%%%%%German%%%%%%%%%%%%%%%%%%%%%%%%%%%%%%
\section*{Declaration of Authorship}
Ich erkläre hiermit gemäß § 9 Abs. 12 APO, dass ich die vorstehende Seminararbeit selbstständig verfasst und keine anderen als die angegebenen Quellen und Hilfsmittel benutzt habe. Des Weiteren erkläre ich, dass die digitale Fassung der gedruckten Ausfertigung der Seminararbeit ausnahmslos in Inhalt und Wortlaut entspricht und zur Kenntnis genommen wurde, dass diese digitale Fassung einer durch Software unterstützten, anonymisierten Prüfung auf Plagiate unterzogen werden kann.\\
\vspace{2\baselineskip}
  
Bamberg, \today

\rule[0.5em]{14em}{0.5pt} \hspace{0.25\linewidth}\rule[0.5em]{14em}{0.5pt}
\vspace{1em}
\hspace{4em} (Ort, Datum) \hspace{0.51\linewidth} (Unterschrift)
}



%%%%%%%%%%%%%%%%%%%%%%%%%%%%%%%%%%%%%%%%%%%%%%%%%%%%%%%%%%%%


%%%%%%%%%%%%%%%%%%%%%%%%%%%%%%%%%%%%%%%%%%%%%%%%%%%%%%%%%%%%


\appendix


\section{Appendix}


Optionally include extra information (complete proofs, additional experiments and plots) in the appendix.
This section will often be part of the supplemental material.


\end{document}
